% !TEX root = ../main.tex

%%% ============================================================
\section{Thực nghiệm và kết quả}
%%% ============================================================

\subsection{Thiết lập \& Câu hỏi nghiên cứu}
\begin{frame}{Câu hỏi nghiên cứu và thiết lập thực nghiệm}
  \begin{columns}[t]
    \begin{column}{0.485\textwidth}
      \uncover<1->{
        \begin{block}{3 Câu hỏi nghiên cứu (RQ)}
          \begin{itemize}\scriptsize\setlength{\itemsep}{5pt}
            \item \textbf{CH1.} Khả năng tự động tạo tác vụ có ý nghĩa và quy mô tập dữ liệu thu được?
            \item \textbf{CH2.} Năng lực của SLM sau fine-tuning so với mô hình nền và các SOTA?
            \item \textbf{CH3.} Hiệu quả chi phí kinh tế của chu trình tự động sinh dữ liệu?
          \end{itemize}
        \end{block}
      }
    \end{column}
    \hfill
    \begin{column}{0.485\textwidth}
      \uncover<2->{
        \begin{block}{Thiết lập thực nghiệm}
          \begin{itemize}\scriptsize\setlength{\itemsep}{3.5pt}
            \item \textbf{Mô hình}: Base Qwen3.5-9B $\rightarrow$ HSGE-9B.
            \item \textbf{Fine-tuning}: LoRA / QLoRA (3 epoch, 1119 bước).
            \item \textbf{Đánh giá}: WebArena 655 task / 4 site, chấm bằng Task SR benchmark.
            \item \textbf{Đối chứng}: cùng prompt, cùng ngân sách bước với mô hình nền.
          \end{itemize}
        \end{block}
      }
    \end{column}

  \end{columns}
\end{frame}

\subsection{CH1: Tự động tạo tác vụ \& Quy mô dữ liệu}
\begin{frame}{CH1 - Tự động tạo tác vụ có ý nghĩa và quy mô dữ liệu}

  \begin{columns}[t,onlytextwidth]
    \begin{column}{0.53\textwidth}
      \vspace{-0.15cm}
      \setlength{\tabcolsep}{3pt}
      \begin{table}
        \scriptsize
        \begin{tabular}{@{}L{3.9cm} R{2.3cm}@{}}
          \toprule
          Quy mô khám phá & \textbf{4 site / 455 URL} \\
          Trajectory ghi nhận & \textbf{3.981 lượt thử} \\
          Trajectory hợp lệ & \textbf{2.117} (53,2\%) \\
          \midrule
          Tri thức tích lũy & \textbf{146 skills / 57 op} \\
          Semantic task & \textbf{142 task} \\
          \midrule
          Dữ liệu huấn luyện & \kw{1.760 trajectory} \\
          Dataset (Train / Val) & \kw{9.007 samples} (8.554 / 453) \\
          \bottomrule
        \end{tabular}
        \vspace{0.04cm}
        \centering
        \\ \scriptsize \textbf{Bảng 1: Quy mô dữ liệu sinh tự động HSGE}
      \end{table}
    \end{column}
    \hfill
    \begin{column}{0.45\textwidth}
      \vspace{-0.15cm}
      \begin{block}{Cổng kiểm soát chất lượng}
        \begin{itemize}\scriptsize\setlength{\itemsep}{0.5pt}
          \item \textbf{LLM-as-Judge}: Loại \kw{801 lượt (27,4\%)} dương tính giả.
          \item \textbf{Effective Action Gate}: Loại \kw{104 trajectory} DOM rỗng.
        \end{itemize}
      \end{block}

      \vspace{-0.25cm}
      \begin{block}{Kết quả thu hoạch}
        \begin{itemize}\scriptsize\setlength{\itemsep}{0.5pt}
          \item Sinh \kw{142 task} từ chuỗi operator.
          \item Thu hoạch \kw{1.760 trajectory sạch} (9.007 mẫu).
        \end{itemize}
      \end{block}
    \end{column}
  \end{columns}

  \vspace{-0.05cm}
  \begin{beamercolorbox}[wd=\textwidth,sep=2.5pt,leftskip=6pt,rightskip=6pt,rounded=true]{block body}
    \scriptsize \kw{Kết luận CH1}: HSGE tự động sinh tác vụ có ý nghĩa và thu hoạch 9.007 mẫu SFT sạch hoàn toàn không cần gán nhãn thủ công.
  \end{beamercolorbox}
\end{frame}

\subsection{CH2: Năng lực của SLM sau fine-tuning}
\begin{frame}{CH2 - So sánh hiệu năng với các hệ thống SOTA}
  \centering
  \vspace{-0.1cm}
  \setlength{\tabcolsep}{4pt}
  \begin{table}
    \scriptsize
    \begin{tabular}{@{}L{4.8cm} R{1.7cm} R{1.2cm} R{1.2cm} R{1.2cm} R{1.2cm}@{}}
      \toprule
      \textbf{Mô hình} & \multicolumn{5}{c}{\textbf{Tỷ lệ thành công Task SR (\%)}} \\
      \cmidrule(lr){2-6}

      & \textbf{Tổng (4 site)} & \textbf{Shopping} & \textbf{Admin} & \textbf{Reddit} & \textbf{GitLab} \\
      \midrule
      \multicolumn{6}{@{}l}{\textit{Mô hình đóng}} \\
      GPT-4o-mini [Gandhi et al.] & 20{,}0 & 19{,}3 & 19{,}2 & 21{,}1 & 20{,}9 \\
      GPT-4o [Gandhi et al.] & 37{,}5 & 32{,}3 & 35{,}7 & 50{,}9 & 36{,}7 \\
      \midrule
      \multicolumn{6}{@{}l}{\textit{Mô hình mở, cỡ nhỏ ($<10$B)}} \\
      Qwen-2.5-7B-Instruct  & 8{,}3 & 9{,}4 & 7{,}1 & 7{,}9 & 8{,}7 \\
      NNetNav-7B [Murty et al.] & 19{,}1 & 20{,}3 & 14{,}3 & 23{,}7 & \textbf{19{,}9} \\
      Go-Browse-7B [Gandhi et al.] & 22{,}6 & 22{,}4 & \textbf{25{,}3} & 30{,}7 & 15{,}3 \\
      \midrule
      Qwen3.5-9B (Mô hình nền) & 13{,}3 & 19{,}8 & 12{,}6 & 6{,}6 & 11{,}1 \\
      HSGE-7B (Qwen-2.5-7B FT) & 16{,}5 & 21{,}4 & 11{,}5 & 23{,}6 & 12{,}2 \\
      \textbf{HSGE-9B (Khóa luận)} & \textbf{24{,}6} & \textbf{28{,}3} & 21{,}4 & \textbf{34{,}0} & 18{,}3 \\
      \bottomrule
    \end{tabular}
    \vspace{0.08cm}
    \centering
    \\ \scriptsize \textbf{Bảng 2: So sánh tỷ lệ thành công Task SR (\%) với các hệ thống SOTA trên WebArena}
  \end{table}
\end{frame}

\begin{frame}{CH2 - So sánh mô hình nền (Base) và fine-tuning theo Website}
  \centering
  \vspace{-0.25cm}
  \setlength{\tabcolsep}{3pt}
  \begin{table}
    \scriptsize
    \begin{tabular}{@{}L{2.3cm} R{2.4cm} R{2.7cm} R{1.9cm} R{2.1cm}@{}}
      \toprule
      \textbf{Website} & \textbf{Base (SR, \%)} & \textbf{Fine-tuned (SR, \%)} &
      \textbf{$\Delta$ SR (đpt)} & \textbf{$p$ (McNemar)} \\
      \midrule
      Shopping       & 37/187 $=$ 19,8 & 53/187 $=$ 28,3 & $+$8,6  & 0,011 \\
      Reddit         & 7/106 $=$ 6,6   & 36/106 $=$ 34,0 & $+$27,4 & $<$0,001 \\
      GitLab         & 20/180 $=$ 11,1 & 33/180 $=$ 18,3 & $+$7,2  & 0,011 \\
      Shopping Admin & 23/182 $=$ 12,6 & 39/182 $=$ 21,4 & $+$8,8  & 0,005 \\
      \midrule
      \textbf{Tổng} & \textbf{87/655 $=$ 13,3} & \textbf{161/655 $=$ 24,6} &
      \kw{$+$11,3} & \kw{$<$0,001} \\
      \bottomrule
    \end{tabular}
    \vspace{0.02cm}
    \centering
    \\ \scriptsize \textbf{Bảng 3: Đối chứng mô hình nền Qwen3.5-9B và mô hình sau Fine-tune theo 4 Website}
  \end{table}

  \vspace{-0.45cm}
  \begin{columns}[t,onlytextwidth]
    \begin{column}{0.49\textwidth}
      \begin{alertblock}{Phân tích kết quả}
        \begin{itemize}\scriptsize\setlength{\itemsep}{0pt}
          \item \textbf{Cú pháp}: Lỗi định dạng giảm từ 26,0\% $\rightarrow$ \kw{0,6\%}.
          \item \textbf{Ra quyết định}: SR chuẩn cú pháp tăng \kw{17,9\% $\rightarrow$ 23,5\%} ($p=0{,}003$).
        \end{itemize}
      \end{alertblock}
    \end{column}
    \hfill
    \begin{column}{0.49\textwidth}
      \begin{block}{Ý nghĩa thống kê McNemar}
        \begin{itemize}\scriptsize\setlength{\itemsep}{0pt}
          \item \textbf{Vượt trội Base}: \kw{100 task độc quyền} HSGE giải được (Base chỉ 26 task).
          \item \textbf{Ý nghĩa thống kê}: Đạt \kw{$p < 0{,}001$} trên toàn bộ 4 website.
        \end{itemize}
      \end{block}
    \end{column}
  \end{columns}

  \vspace{-0.15cm}
  \begin{beamercolorbox}[wd=\textwidth,sep=2.5pt,leftskip=6pt,rightskip=6pt,rounded=true]{block body}
    \scriptsize \kw{Kết luận CH2}: Fine-tuning nâng Task SR lên 24,6\% ($+11,3\%$, $p < 0{,}001$), làm chủ cả cú pháp lẫn ra quyết định, đạt mức SOTA $<10$B.
  \end{beamercolorbox}
\end{frame}

\subsection{CH3: Hiệu quả chi phí trong tự động tạo dữ liệu}
\begin{frame}{CH3 - Hiệu quả chi phí trong quá trình tự động tạo dữ liệu}
  \vspace{-0.15cm}
  \begin{columns}[t,onlytextwidth]
    \begin{column}{0.49\textwidth}
      \begin{block}{Chi phí API toàn chu trình \& Token}
        \begin{itemize}\scriptsize\setlength{\itemsep}{2pt}
          \item \textbf{OpenAI API}: \$42,29 (\texttt{solver}, \texttt{explorer}).
          \item \textbf{DeepSeek API}: \$11,92 (\texttt{judge}, pipeline skill).
          \item \textbf{Tổng chi phí API}: \kw{\$54,21} ($<$ \$14 / site).
          \item \textbf{Đơn giá / trajectory}: \kw{\$0,0308} / trajectory.
          \item \textbf{Lời gọi API}: \kw{20.808 lượt} LLM (Solver \kw{85,5\%}).
          \item \textbf{Tối ưu token}: Chuyển sang GPT-5.6-luna giúp giảm đơn giá token đầu vào \kw{$\sim$12 lần} so với GPT-4o.
        \end{itemize}
      \end{block}
    \end{column}
    \hfill
    \begin{column}{0.48\textwidth}
      \begin{block}{So sánh hiệu quả chi phí với Go-Browse}
        \centering
        \vspace{-0.1cm}
        \setlength{\tabcolsep}{3pt}
        \scriptsize
        \begin{tabular}{@{}L{2.3cm} R{3.3cm}@{}}
          \toprule
          \textbf{Phương pháp} & \textbf{Chi phí API / Trajectory} \\
          \midrule
          Go-Browse-7B & \$0,1026 \\
          \textbf{HSGE (Khóa luận)} & \kw{\$0,0308} \\
          \midrule
          \textbf{Tỷ lệ tiết kiệm} & \kw{Rẻ hơn 3,3 lần} \\
          \bottomrule
        \end{tabular}

        \vspace{0.02cm}
        \scriptsize \textbf{Bảng 4: So sánh đơn giá API trajectory tự động tạo}
      \end{block}

      \vspace{-0.25cm}
      \begin{block}{Ý nghĩa kinh tế}
        \begin{itemize}\scriptsize\setlength{\itemsep}{1pt}
          \item Khẳng định tính khả thi với tài nguyên giới hạn.
          \item Chu trình HSGE nằm trong ngân sách nhóm nhỏ.
        \end{itemize}
      \end{block}
    \end{column}
  \end{columns}

  \vspace{0.05cm}
  \begin{beamercolorbox}[wd=\textwidth,sep=2.5pt,leftskip=6pt,rightskip=6pt,rounded=true]{block body}
    \scriptsize \kw{Kết luận CH3}: Chi phí API tự tạo dữ liệu \$0,0308/trajectory (rẻ hơn Go-Browse 3,3 lần), chứng minh tính khả thi kinh tế vượt trội cho nhóm nhỏ.
  \end{beamercolorbox}
\end{frame}
